The preceding sections separate real zero curves from complex roots because they
have different codimensions.  A polynomial family can nevertheless carry both
objects at once.  Its imaginary part supplies a real zero carrier, while its full
complex zero divisor supplies a finite thread inside that carrier.  This is the
first construction in the paper in which the relation between slices, rather
than only the geometry of one slice, is encoded by the same polynomial that
selects the thread.

\subsection{A polynomial defines both carrier and thread}

Let $B$ be an oriented one-manifold, let $D\subset\C$ be a disc, and let
\begin{equation}
 P:D\times B\longrightarrow\C
 \label{eq:polynomial-carrier-field}
\end{equation}
be smooth in the parameter and holomorphic in the spatial variable.  Put
\begin{equation}
 A=\im P,
 \qquad
 \Sigma_P=A^{-1}(0),
 \qquad
 K_P=P^{-1}(0).
 \label{eq:polynomial-carrier-thread}
\end{equation}
Then $K_P\subset\Sigma_P$ tautologically, but the two projections to $B$ have
different regularity conditions.

\begin{proposition}[Polynomial carrier and intrinsic root thread]
\label{prop:polynomial-carrier-thread}
Assume that every root of $P_b$ is simple, remains in the interior of $D$, and
that the number of roots is constant.  Then $K_P\to B$ is a proper finite
covering.  Suppose in addition that $\dd A$ is nonzero on $\Sigma_P$.  Then
$\Sigma_P$ is a smooth cooriented surface, and the critical points of
$\Sigma_P\to B$ are exactly the points satisfying
\begin{equation}
 \partial_zP_b(z)=0,
 \qquad
 \im P_b(z)=0.
 \label{eq:carrier-critical-value-real}
\end{equation}
In particular, the root thread avoids every such critical point.

On the complement of the spatial critical set of $P_b$, for every $\mu\ne0$
and every real $\lambda$, the metric
\begin{equation}
 g_b=
 \frac{|\partial_zP_b|^2}
      {\mu^2+\lambda^2(\im P_b)^2}|\dd z|^2
 \label{eq:polynomial-carrier-aeg-metric}
\end{equation}
satisfies
$|\nabla A_b|_{g_b}^2=\mu^2+\lambda^2A_b^2$.
Thus a spatial critical point may be declared as an essential singular stratum,
while the roots remain regular finite threads.
\end{proposition}

\begin{proof}
At a root, simplicity says $\partial_zP_b\ne0$.  The real implicit-function
theorem therefore writes $K_P$ locally as disjoint graphs over $B$.  The fixed
finite degree and interior control make this covering proper.

The hypothesis $\dd A\ne0$ gives the smooth cooriented carrier.  Its projection
to $B$ is critical precisely when the vertical derivative of $A$ vanishes.
Because $P_b$ is holomorphic, this vertical real one-form vanishes precisely
when $\partial_zP_b=0$, proving \eqref{eq:carrier-critical-value-real}.  Such a
point cannot be a root under the square-free hypothesis.  Finally,
the Euclidean norm squared of $\dd\im P_b$ is
$|\partial_zP_b|^2$.  Conformal inversion in
\eqref{eq:polynomial-carrier-aeg-metric} proves the AEG equation.
\end{proof}

The proposition distinguishes two discriminants.  A collision of roots belongs
to the algebraic discriminant of $P$.  A Morse carrier reconnection occurs
instead when a nondegenerate critical value of a still square-free polynomial
crosses the real axis transversely.
The thread can therefore remain smooth while the real carrier passes through a
Morse event.

There is also a canonical geometric direction along the thread once the spatial
complex coordinate has been declared.

\begin{proposition}[Carrier framing and discriminant period]
\label{prop:carrier-framing-discriminant}
Under the hypotheses of Proposition~\ref{prop:polynomial-carrier-thread}, assume
in addition that $B=S^1$ and every $P_b$ is a monic polynomial of one fixed
degree $n$.  Write
the roots locally as $r_1(b),\ldots,r_n(b)$ and define
\begin{equation}
 v_i(b)=\frac{1}{\partial_zP_b(r_i(b))}\in T_{r_i(b)}D\cong\C.
 \label{eq:polynomial-carrier-framing-vector}
\end{equation}
Then $v_i$ is a nonzero vertical tangent vector to $\Sigma_P$ along the root
thread.  It therefore frames the thread inside the carrier.  Around a parameter
circle, the sum of its rotation numbers over the root cycles is
\begin{equation}
 \operatorname{Fr}_{\Sigma_P}(K_P)
 =-\operatorname{wind}\!\left(\operatorname{disc}_zP\right).
 \label{eq:carrier-framing-discriminant-period}
\end{equation}
The integer is relative to the declared complex trivialization of the spatial
disc; in a nontrivial surface bundle it must be replaced by the corresponding
relative Euler or framing class.
\end{proposition}

\begin{proof}
At a root, $\dd P_b(v_i)=1$.  Its imaginary part vanishes, so
$\dd A_b(v_i)=0$ and $v_i$ lies in the vertical tangent line of the carrier.
For a monic polynomial,
\begin{equation}
 \prod_{i=1}^n\partial_zP_b(r_i)
 =(-1)^{n(n-1)/2}\operatorname{disc}_zP_b.
 \label{eq:derivative-product-discriminant}
\end{equation}
The constant sign has zero winding.  Taking inverses and summing the winding
over every root cycle proves \eqref{eq:carrier-framing-discriminant-period}.
\end{proof}

\subsection{The quadratic threaded-carrier family}

Fix an integer $m\ge1$ and a radius $R>1$.  Let
$D_R=\{z\in\C:|z|\le R\}$ be the closed disc, and write
$t=e^{i\theta}\in S^1$ and $z=x+iy\in D_R$.  Consider
\begin{align}
 P_m(z,t)&=z^2-t^m,
 \label{eq:quadratic-thread-polynomial}\\
 A_m(x,y,\theta)&=\im P_m
                 =2xy-\sin(m\theta),
 \label{eq:quadratic-thread-assignment}\\
 \Sigma_m&=A_m^{-1}(0)\subset D_R\times S^1,
 \qquad
 K_m=P_m^{-1}(0)\subset\Sigma_m.
 \label{eq:quadratic-thread-carrier}
\end{align}
The parameter $t$ is a deformation register.  It is neither the AEG strength
$\lambda$ nor the fixed Hecke translation length $\sqrt2$.

\begin{theorem}[Quadratic threaded carriers]
\label{thm:quadratic-threaded-carriers}
For every $m\ge1$ and $R>1$, the following statements hold.
\begin{enumerate}
\item The carrier $\Sigma_m$ is a smooth, compact, connected, cooriented surface
  meeting $\partial D_R\times S^1$ neatly.  Its boundary has four connected
  components.
\item The projection $p_m:\Sigma_m\to S^1$ has exactly $2m$ critical points,
  all Morse saddles.  They occur at
  \begin{equation}
   z=0,
   \qquad
   \theta=\frac{k\pi}{m},
   \qquad 0\le k<2m.
   \label{eq:quadratic-carrier-saddles}
  \end{equation}
  Consequently
  \begin{equation}
   \chi(\Sigma_m)=-2m,
   \qquad
   \Sigma_m\cong\Sigma_{g,4},
   \qquad
   g=m-1,
   \label{eq:quadratic-carrier-genus}
  \end{equation}
  where $\Sigma_{g,4}$ denotes an oriented genus-$g$ surface with four boundary
  components.
\item On every punctured slice $D_R\setminus\{0\}$,
  \begin{equation}
   g_{m,\theta}
   =\frac{4|z|^2}
          {\mu^2+\lambda^2A_m(z,\theta)^2}|\dd z|^2
   \label{eq:quadratic-carrier-singular-aeg}
  \end{equation}
  is an exact AEG metric.  The center is an essential metric singularity whose
  metric completion has cone angle $4\pi$.  At
  the parameters in \eqref{eq:quadratic-carrier-saddles}, the slice zero germ is
  the four-pronged germ $2xy=0$.
\item The root set $K_m\to S^1$ is a smooth two-sheeted finite thread disjoint
  from all carrier saddles.  As a closed two-braid in $D_R\times S^1$, it is
  \begin{equation}
   \beta_m=\sigma_1^m,
   \qquad
   \widehat\beta_m=T(2,m).
   \label{eq:quadratic-carrier-braid}
  \end{equation}
  It has $\gcd(2,m)$ connected components.
\end{enumerate}
\end{theorem}

\begin{proof}
The total gradient of $A_m$ is
\begin{equation}
 \dd A_m=(2y,2x,-m\cos(m\theta)).
 \label{eq:quadratic-carrier-total-gradient}
\end{equation}
If its first two entries vanish at a point of $\Sigma_m$, then
$\sin(m\theta)=0$ and the last entry is nonzero.  Hence $\Sigma_m$ is smooth;
it is compact as a closed subset of $D_R\times S^1$, and it is cooriented by
$\dd A_m$.

Failure of neatness at $|z|=R$ would require the spatial part
$(2y,2x)$ to be parallel to the radial vector $(x,y)$ and also
$\cos(m\theta)=0$.  Parallelism forces $y=x$ or $y=-x$, while the carrier
equation then forces $|z|=1$.  This contradicts $R>1$.  On the boundary write
$z=Re^{i\phi}$.  The equation is
\begin{equation}
 R^2\sin(2\phi)=\sin(m\theta).
 \label{eq:quadratic-carrier-boundary-equation}
\end{equation}
Since $R^{-2}<1$, its four solutions are four disjoint periodic graphs over
$S^1$.  Thus the boundary has four components.

The projection is critical exactly where the spatial derivative of $A_m$
vanishes, giving \eqref{eq:quadratic-carrier-saddles}.  Near such a point, put
$\theta=\theta_0+s$.  Solving the carrier equation for $s$ gives
\begin{equation}
 s=(-1)^k\frac{2}{m}xy+O\bigl((xy)^3\bigr).
 \label{eq:quadratic-carrier-saddle-normal-form}
\end{equation}
The Hessian has one positive and one negative eigenvalue, so every critical
point is an index-one saddle.

At $\theta=0$, the slice zero is the union of the two coordinate diameters.
Its four boundary points are joined at the central point, so all four boundary
circles of the total carrier belong to one connected component.  Every regular
fiber component is an interval ending on the boundary, and every critical point
is one of the central saddles.  Hence every point of $\Sigma_m$ lies in that
same component.

Cut $\Sigma_m$ along a regular fiber $F$, which is the disjoint union of two
closed intervals.  The product $F\times[0,1]$ has Euler characteristic two, and
the circle-valued Morse movie attaches one one-handle for each of the $2m$
saddles, giving Euler characteristic $2-2m$ before regluing.  Identifying the
two end copies of $F$ subtracts $\chi(F)=2$.  Hence
$\chi(\Sigma_m)=-2m$.  With four boundary components,
$2-2g-4=-2m$, proving \eqref{eq:quadratic-carrier-genus}.

The metric claim follows from Proposition~\ref{prop:polynomial-carrier-thread},
since $|\partial_zP_m|^2=4|z|^2$.  The coefficient tends to zero at the center,
so it cannot extend there as a positive-definite metric.  In the metric
completion the substitution $\rho=C_\theta |z|^2/2+O(|z|^4)$ gives angular
term $4\rho^2\dd\arg(z)^2$, hence cone angle $4\pi$.  When
$\sin(m\theta)=0$, the zero equation is $2xy=0$.

Finally, the two roots are
\begin{equation}
 z_\pm(\theta)=\pm e^{im\theta/2}.
 \label{eq:quadratic-carrier-roots}
\end{equation}
They lie on $|z|=1$, avoid the center, and determine the standard closed
two-braid $\sigma_1^m$.  Its closure is $T(2,m)$ and has
$\gcd(2,m)$ components.
\end{proof}

Thus this single quadratic template produces carriers of every nonnegative
genus, with a fixed four-component boundary.  It is not universal among all
surfaces or all links: its threads are precisely the two-strand torus links.
Its value is that the carrier reconnections and the thread are forced by one
displayed polynomial rather than selected independently.
Although $\Sigma_m\to S^1$ is proper and its total incidence is smooth and
neat, it is not a proper zero tube in Definition~\ref{def:zero-object-hierarchy}:
its $2m$ saddles are precisely the failures of the submersion condition.

\subsection{Arithmetic components and the horizontal exact sequence}

Let $k$ be a field of characteristic different from two and set
$B=\mathbb G_{m,k}$.  The same formula defines a relative divisor
\begin{equation}
 \mathcal D_m
 =V(u^2-t^m)\subset\mathbb A^1_u\times B.
 \label{eq:quadratic-carrier-relative-divisor}
\end{equation}

\begin{proposition}[Arithmetic components of the thread]
\label{prop:quadratic-thread-arithmetic-components}
The projection $\mathcal D_m\to B$ is finite \'etale of degree two.  The polynomial
$u^2-t^m$ is irreducible over $k(t)$ if and only if $m$ is odd.  If $m$ is even,
then
\begin{equation}
 u^2-t^m=(u-t^{m/2})(u+t^{m/2}).
 \label{eq:quadratic-thread-factorization}
\end{equation}
Let $\operatorname{Irr}_{k(t)}(\mathcal D_m)$ denote the irreducible components
of the generic fiber.  Consequently, in this family the following three numbers
agree:
\begin{equation}
 \#\operatorname{Irr}_{k(t)}(\mathcal D_m)
 =\#\pi_0(K_m)
 =\#\pi_0\bigl(T(2,m)\bigr)
 =\gcd(2,m).
 \label{eq:quadratic-thread-component-identity}
\end{equation}
\end{proposition}

\begin{proof}
The polynomial is monic and its derivative is $2u$.  A common zero of the
polynomial and its derivative would require $u=t=0$, which is impossible on
$B$.  Thus the cover is finite \'etale.  If $m$ is odd, the $t$-adic valuation of
$t^m$ is odd, so it is not a square in $k(t)$ and the quadratic is irreducible.
For even $m$, \eqref{eq:quadratic-thread-factorization} proves the converse.
The remaining equalities follow from Theorem~\ref{thm:quadratic-threaded-carriers}.
\end{proof}

The full horizontal information is finer than the component count.  The
discriminant in $u$ is
\begin{equation}
 \Delta_m(t)=\operatorname{disc}_u(u^2-t^m)=4t^m.
 \label{eq:quadratic-thread-discriminant}
\end{equation}
It never vanishes on $B$, so the roots never collide.  On the parameter circle,
however, the $2m$ carrier saddles occur exactly when
$\im\Delta_m(t)=0$.  They are real critical-value wall crossings, not points of
the algebraic root discriminant.
The standard exact sequence
\begin{equation}
 1\longrightarrow P_2\longrightarrow B_2
 \longrightarrow S_2\longrightarrow1
 \label{eq:two-braid-forgetting-exact-sequence}
\end{equation}
records the loss from a two-strand braid to its root permutation.  Under the
identifications
\begin{equation}
 P_2\cong\Z\langle\sigma_1^2\rangle,
 \qquad
 B_2\cong\Z\langle\sigma_1\rangle,
 \qquad
 S_2\cong\Z/2\Z,
 \label{eq:two-braid-integer-identifications}
\end{equation}
it becomes
\begin{equation}
 0\longrightarrow\Z\xrightarrow{\,\times2\,}\Z
 \xrightarrow{\bmod 2}\Z/2\Z\longrightarrow0.
 \label{eq:two-braid-integer-exact-sequence}
\end{equation}

\begin{theorem}[The threaded-carrier identity]
\label{thm:threaded-carrier-identity}
With the positive orientation conventions above,
\begin{equation}
 \boxed{
 \operatorname{ord}_{t=0}\Delta_m
 =\frac{1}{2\pi i}\int_{|t|=1}\dd\log\Delta_m
 =\operatorname{exp}(\beta_m)
 =-\operatorname{Fr}_{\Sigma_m}(K_m)
 =-\frac12\chi(\Sigma_m)
 =g(\Sigma_m)+1
 =m.}
 \label{eq:threaded-carrier-main-identity}
\end{equation}
The image of this integer in $S_2$ is $m\bmod2$, which is exactly the
permutation monodromy and exactly the parity governing
\eqref{eq:quadratic-thread-component-identity}.  If $m$ is even, then
$\beta_m\in P_2$ and
\begin{equation}
 \operatorname{lk}(K_m^+,K_m^-)
 =\frac12\operatorname{exp}(\beta_m)
 =\frac m2,
 \label{eq:quadratic-thread-linking}
\end{equation}
where the sign changes under orientation reversal.
\end{theorem}

\begin{proof}
Equation \eqref{eq:quadratic-thread-discriminant} gives both its order at the
missing point $t=0$ of the affine compactification and its logarithmic period.
For two unordered roots, the discriminant map to $\C^\times$ sends the positive
half-twist $\sigma_1$ to a loop of winding one.  It therefore identifies
$B_2$ with $\pi_1(\C^\times)$, so the logarithmic period is the braid exponent.
The surface equalities are \eqref{eq:quadratic-carrier-genus}.  Reduction modulo
two is the right-hand map in \eqref{eq:two-braid-integer-exact-sequence}.
Finally, the closure of $\sigma_1^{2r}$ has two components with linking number
$r$, giving \eqref{eq:quadratic-thread-linking}.
\end{proof}

This theorem is a horizontal, parameter-dependent statement.  Arithmetic
factorization and permutation monodromy retain only the parity layer.  The
analytic logarithmic period retains the full braid exponent; on the pure-braid
kernel, half of that period---equivalently, its coordinate relative to the
generator $\sigma_1^2$---is mutual linking.  No cohomological transgression from
Paper~I's
slice torsion has been asserted.  Over a circle base, the relevant proved
structure is the monodromy exact sequence
\eqref{eq:two-braid-forgetting-exact-sequence}, not a Serre transgression.

\subsection{The q=4 toric polynomial normal form}

Theorem~\ref{thm:q4-sign-cover-torus-link} computed more than the abstract link
type.  On the lifted median Heegaard torus, with the two filling meridians denoted
$m_P,m_Q$, each component of the lifted cusp link has primitive slope
\begin{equation}
 a_\infty=2m_P+m_Q.
 \label{eq:q4-thread-slope-recall}
\end{equation}
The marked slope and deck data select a binomial normal form.

\begin{corollary}[q=4 cusp link as the polynomial thread]
\label{cor:q4-polynomial-thread}
Choose angular coordinates $(u,t)\in S^1\times S^1$ whose oriented circle
classes are $(m_P,m_Q)$.  Then the two lifted cusp components are isotopic on
the marked torus to
\begin{equation}
 \{u^2=t^4\}
 =\{u=t^2\}\amalg\{u=-t^2\}.
 \label{eq:q4-toric-binomial}
\end{equation}
Consequently, after choosing a meridional radial coordinate, pushing the two
boundary curves slightly into the filling solid torus, and normalizing their
radius to one in a closed disc $D_R$ with $R>1$, the resulting closed braid is
the thread $K_4$ of Theorem~\ref{thm:quadratic-threaded-carriers}.  For this
additional radial carrier choice,
\begin{equation}
 \begin{gathered}
 \Sigma_4\cong\Sigma_{3,4},
 \qquad
 \#\operatorname{Crit}(p_4)=8,\\
 \operatorname{wind}(\Delta_4)=4,
 \qquad
 \beta_4=\sigma_1^4\in P_2,
 \qquad
 \operatorname{lk}(K_4^+,K_4^-)=2.
 \end{gathered}
 \label{eq:q4-threaded-carrier-data}
\end{equation}
The root permutation is trivial even though its pure-braid class is nonzero.
\end{corollary}

\begin{proof}
In the notation of the slope computation in
Theorem~\ref{thm:q4-sign-cover-torus-link}, the covering lattice is
\(
 \Lambda=\{(x,y)\in\Z^2:x+y\equiv0\pmod2\}
\), and $(m_P,m_Q)$ is a basis of $\Lambda$.  A nontrivial deck translation is
represented by $(1,0)$ and hence, modulo $\Lambda$, by
$(m_P+m_Q)/2$.  It acts on the dual torus characters as
\begin{equation}
 (u,t)\longmapsto(-u,-t).
 \label{eq:q4-toric-deck-action}
\end{equation}
The primitive character annihilating the slope $2m_P+m_Q$ is
$r=ut^{-2}$.  Thus one cusp lift is $r=c$, for some unit phase $c$, and the
deck-related lift is $r=-c$.  Their union is
$r^2=c^2$, or $u^2=c^2t^4$.  Rescaling the $u$ character gives
\eqref{eq:q4-toric-binomial}.  Equivalently, the curves $u=\pm t^2$ each have
class $2m_P+m_Q$ on the marked torus.

The remaining statements are the case $m=4$ of
Theorems~\ref{thm:quadratic-threaded-carriers} and
\ref{thm:threaded-carrier-identity}.
\end{proof}

The binomial can in fact be recovered without first choosing a radial disc
extension.  The intrinsic extension is the affine cone of the logarithmic tangent
line of the compactified sign-cover orbifold.

\begin{theorem}[The q=4 logarithmic tangent cone]
\label{thm:q4-logarithmic-tangent-cone}
Let $X$ be the compact complex orbifold obtained from
$\mathcal O_4^\chi$ by adjoining its two cusp points, equipped with the
inherited normalized Hauptmodul $\beta$ and deck involution $\iota$, and let
$D_\infty\subset X$ be their reduced sum.  Then the following statements hold.
\begin{enumerate}
\item As a marked orbifold pair,
  \begin{equation}
   (X,D_\infty)
   \cong
   \bigl(\mathbb P(2,1),V(U^2-V^4)\bigr),
   \label{eq:q4-weighted-orbifold-pair}
  \end{equation}
  where $U,V$ have weights two and one.  In these weighted coordinates,
  \begin{equation}
   \beta=\frac{U^2}{U^2-V^4}.
   \label{eq:q4-weighted-hauptmodul}
  \end{equation}
  The residual order-two point is $[1:0]$, the regular point lying over the old
  order-two elliptic point is $[0:1]$, and $\iota[U:V]=[-U:V]$.
\item The logarithmic tangent line is
  \begin{equation}
   L:=T_X(-\log D_\infty)\cong\mathcal O_{\mathbb P(2,1)}(-1).
   \label{eq:q4-log-tangent-line}
  \end{equation}
  Its canonical section ring is
  \begin{equation}
   R_L=\bigoplus_{d\ge0}H^0(X,L^{-d})
   \cong\C[V,U],
   \qquad \deg V=1,\quad\deg U=2,
   \label{eq:q4-log-tangent-section-ring}
  \end{equation}
  and the section cutting out $D_\infty$ is, up to a nonzero scalar,
  $U^2-V^4$.  Thus the graded affine cone of the marked pair canonically
  recovers the binomial divisor as a graded-isomorphism class.
\item The standard hyperbolic cusp compactification of
  $T^1\mathcal O_4^\chi$ is fiber-preservingly diffeomorphic to the weighted
  Hopf circle bundle
  \begin{equation}
   S^3\longrightarrow\mathbb P(2,1),
   \qquad
   e^{i\alpha}(U,V)=(e^{2i\alpha}U,e^{i\alpha}V).
   \label{eq:q4-weighted-hopf-fibration}
  \end{equation}
  Under this diffeomorphism the two filled cusp fibers are exactly
  \begin{equation}
   S^3\cap V(U^2-V^4)=T(2,4).
   \label{eq:q4-weighted-polynomial-link}
  \end{equation}
  The deck involution lifts to the antipodal map
  $(U,V)\mapsto(-U,-V)$, which exchanges the two link components and has
  quotient $\mathbb {RP}^3$.
\end{enumerate}
\end{theorem}

\begin{proof}
Put
\(
 r=(1+W)/(1-W)=-iF
\).
Since $F^2=\beta/(1-\beta)$,
\begin{equation}
 r^2=\frac{\beta}{\beta-1},
 \qquad
 \beta=\frac{r^2}{r^2-1}.
 \label{eq:q4-r-beta-coordinate}
\end{equation}
Proposition~\ref{prop:q4-sign-cover-cylinder} shows that $r$ is a coarse
coordinate on the compactified sign cover.  The residual order-two point is
$r=\infty$, while the two cusps are $r=1$ and $r=-1$.  The weighted projective
line
\begin{equation}
 \mathbb P(2,1)
 =[(\C^2\setminus\{0\})/\C^\times],
 \qquad
 c\cdot(U,V)=(c^2U,cV),
 \label{eq:weighted-projective-line-definition}
\end{equation}
has the same coarse coordinate $r=U/V^2$ and the same order-two stabilizer at
$V=0$.  Substitution in \eqref{eq:q4-r-beta-coordinate} gives
\eqref{eq:q4-weighted-hauptmodul}; its polar divisor is
$U^2-V^4=0$.  This proves the marked-pair statement.

The weighted Euler sequence is
\begin{equation}
 0\longrightarrow\mathcal O
 \longrightarrow\mathcal O(2)\oplus\mathcal O(1)
 \longrightarrow T_X\longrightarrow0,
 \label{eq:q4-weighted-euler-sequence}
\end{equation}
so $T_X\cong\mathcal O(3)$.  Each ordinary cusp point is cut out by a section
of $\mathcal O(2)$, and their sum is the divisor of $U^2-V^4$, of class
$\mathcal O(4)$.  On an orbifold curve the logarithmic tangent line is
$T_X(-D_\infty)$; hence it is $\mathcal O(3-4)=\mathcal O(-1)$.  The global
sections of $\mathcal O(d)$ are the weighted homogeneous polynomials of degree
$d$, proving \eqref{eq:q4-log-tangent-section-ring} and the asserted equation
of the canonical divisor section.

The unit circle bundle of $\mathcal O(-1)$ is $S^3$ with the weighted action
\eqref{eq:q4-weighted-hopf-fibration}.  Away from $D_\infty$, the logarithmic
tangent line equals the ordinary tangent line.  Radial normalization from the
hyperbolic norm to any smooth Hermitian norm on the logarithmic line gives a
fiber-preserving identification of their unit-circle bundles over the punctured
orbifold.  Near a cusp coordinate $q=0$, the logarithmic frame is
$q\partial_q$, and the positively tangent horocycle vector is represented by
$i q\partial_q$.  Its phase is therefore constant.  The horocycle-tangent loop
is the boundary of the base disc at that fixed fiber phase when the logarithmic
circle bundle is extended over $q=0$.  This is exactly the cusp-filling convention of
Theorem~\ref{thm:q4-sign-cover-torus-link}, proving the fiber-preserving
compactification statement.

The preimage of $D_\infty$ is
$S^3\cap V(U^2-V^4)$.  Its two branches $U=\pm V^2$ are regular weighted Hopf
fibers and, on their common standard torus, have primitive slope $(2,1)$.
Their union is $T(4,2)=T(2,4)$.

It remains to identify the deck lift rather than merely its action on the base.
On the affine $r$-chart, the logarithmic tangent is generated by
\begin{equation}
 e=(r^2-1)\partial_r.
 \label{eq:q4-log-tangent-deck-frame}
\end{equation}
For $\iota(r)=-r$ one has $\dd\iota(e)=-e$.  Under
$\mathcal O(-1)^\times\cong\C^2\setminus\{0\}$, write
$(U,V)=(\ell^2r,\ell)$.  The derivative lift is
$(r,\ell)\mapsto(-r,-\ell)$, equivalently
$(U,V)\mapsto(-U,-V)$.  It is free, commutes with the weighted circle action,
and exchanges the two branches, so it is the extended unit-tangent deck
involution.  Its quotient is $\mathbb {RP}^3$.
\end{proof}

\begin{remark}[Arithmetic quadratic twist]
\label{rem:q4-log-cone-arithmetic-twist}
Theorem~\ref{thm:q4-logarithmic-tangent-cone} is a complex-orbifold statement.
Over $K=\Q(\sqrt2)$, retaining the square-root coordinate $F=U/V^2$ gives
instead
\begin{equation}
 \beta=\frac{U^2}{U^2+V^4},
 \qquad
 D_\infty=V(U^2+V^4).
 \label{eq:q4-arithmetic-plus-twist}
\end{equation}
The change $r=-iF$ used above converts this to the split minus form only after
adjoining $i$.  The plus form is prime over $K(V)$ and splits over $K(i)$ into
$U=\pm iV^2$, whose two complex link components are exchanged by
$\operatorname{Gal}(K(i)/K)$.  Alternatively, declaring the two geometric cusps
individually $K$-rational supplies a split $K$-descent, which can be written in a
$K$-coordinate $W$ with cusp values $0$ and $\infty$.  This is a different
descent datum from the $F$-model above.  The complex sign cover alone does not
choose between these constant quadratic twists.  A future arithmetic
naturality theorem must declare the descent datum before claiming either
$K$-primality or two $K$-rational cusp components.
\end{remark}

The normal form is canonical only at the correct level.  The two filling
meridians determine the integral torus marking up to the declared orientation
and unit changes, and the slope with its two deck-related components determines
the Laurent binomial divisor up to those gauges.  Extending an angular coordinate
radially to a disc is additional carrier data.  Accordingly,
Corollary~\ref{cor:q4-polynomial-thread} does not identify $\Sigma_4$ with the
two-dimensional Hecke zero dessin, nor does it derive a coefficient family from
an unrestricted marked history.  It proves instead that the global q=4 filling
data force a toric polynomial representative whose thread is the same cusp link.
Theorem~\ref{thm:q4-logarithmic-tangent-cone} strengthens this to a canonical
graded-isomorphism class: the binomial is the cusp-divisor section in the section
ring of the logarithmic tangent line.  Its affine-cone vertex is not the Hecke
order-four point and is not the local four-pronged AEG zero germ.  Passing from
this complex cone to the real carrier $\Sigma_4$ by treating $V/|V|$ as an
external slice parameter remains the radial extension choice just described.

\begin{warning}[The pure-braid and unit-tangent extensions differ]
\label{warn:q4-braid-versus-unit-tangent-extension}
The exact sequence \eqref{eq:two-braid-forgetting-exact-sequence} has quotient
$S_2$, not $\Gamma_4$.  The object that can be compared with a $q=4$ extension
is its pullback
\begin{equation}
 \chi^*B_2
 =\{(\gamma,b)\in\Gamma_4\times B_2:
       \chi(\gamma)=\operatorname{perm}(b)\}.
 \label{eq:q4-pullback-two-braid-extension}
\end{equation}
Write $\Gamma_4=\langle J,R\mid J^2=R^4=1\rangle$.  Any lifts of the two
sign-odd generators to $B_2\cong\Z$ have odd exponents.  In the pure-braid
coordinate where $\sigma_1^2$ is one, $J^2$ therefore contributes an odd
integer and $R^4$ contributes twice an odd integer.  The latter is nonprimitive.
By contrast, the natural elliptic lifts in the unit-tangent extension have
relator powers $\widetilde J^{\,2}=f^{\pm1}$ and
$\widetilde R^{\,4}=f^{\pm1}$, with signs fixed by rotation conventions and
$f$ a primitive tangent-fiber class.  Thus the pullback two-braid extension is
not the unit-tangent extension.  In particular, the $q=4$ linking number cannot
be renamed unit-tangent fiber winding, or Paper~I torsion, without a new
order-weighted framing construction.
\end{warning}

The parity obstruction is specific to two strands.  Four strands carry a
central extension whose two elliptic relators have the same primitive residue.
This gives a sharper exact-sequence calibration, although it still does not
supply the desired history-to-coefficient-path functor.

\begin{proposition}[Four-strand central-extension calibration]
\label{prop:q4-four-strand-central-calibration}
In $B_4$, put
\begin{equation}
 \delta_4=\sigma_1\sigma_2\sigma_3,
 \qquad
 \Delta_4=(\sigma_1\sigma_2\sigma_3)(\sigma_1\sigma_2)\sigma_1,
 \qquad
 z_4=\Delta_4^2=\delta_4^4.
 \label{eq:q4-four-strand-garside-elements}
\end{equation}
Let $\overline B_4=B_4/\langle z_4\rangle$.  With the positive tangent-fiber
class denoted by $f$, choose the standard determinant-one elliptic lifts, whose
tangent rotations are clockwise.  The unit-tangent extension has presentation
\begin{equation}
 \widetilde\Gamma_4
 =\left\langle \widetilde J,\widetilde R,f\ \middle|\
 [f,\widetilde J]=[f,\widetilde R]=1,\
 \widetilde J^2=\widetilde R^4=f^{-1}\right\rangle,
\label{eq:q4-unit-tangent-central-presentation}
\end{equation}
and fits into
$1\to\langle f\rangle\to\widetilde\Gamma_4\to\Gamma_4\to1$.
Then
\begin{equation}
 \psi:\Gamma_4\longrightarrow\overline B_4,
 \qquad
 J\longmapsto[\Delta_4^{-1}],
 \qquad
 R\longmapsto[\delta_4^{-1}]
 \label{eq:q4-four-strand-quotient-map}
\end{equation}
is a homomorphism, and the lift
\begin{equation}
 \widetilde J\longmapsto\Delta_4^{-1},
 \qquad
 \widetilde R\longmapsto\delta_4^{-1},
 \qquad
 f\longmapsto z_4
 \label{eq:q4-four-strand-central-lift}
\end{equation}
identifies the unit-tangent extension with the pullback extension
\begin{equation}
 \widetilde\Gamma_4
 \cong
 \Gamma_4\mathbin{\times}_{\overline B_4}B_4.
 \label{eq:q4-four-strand-pullback-extension}
\end{equation}

Choose a standard four-point root configuration and realize the braid words in
\eqref{eq:q4-four-strand-garside-elements} by square-free monic polynomial
coefficient loops.  Both elliptic relator loops then give the same inverse full
twist $z_4^{-1}$, and
\begin{equation}
 \frac1{12}\frac1{2\pi i}\int\dd\log\operatorname{disc}P_{\widetilde J^2}
 =
 \frac1{12}\frac1{2\pi i}\int\dd\log\operatorname{disc}P_{\widetilde R^4}
 =-1.
 \label{eq:q4-four-strand-normalized-discriminant-period}
\end{equation}
This is their common exponent relative to the positive fiber generator.
\end{proposition}

\begin{proof}
The standard Garside identities in $B_4$ give
$\Delta_4^2=\delta_4^4=z_4$, the positive full twist
\cite[Chapter~1]{Birman1974Braids}.  Therefore the two assignments in
\eqref{eq:q4-four-strand-quotient-map} satisfy $J^2=R^4=1$ after quotienting by
$\langle z_4\rangle$.  With the declared clockwise arguments, the local
derivatives of $J$ and $R=JT$ have rotation angles $-\pi$ and $-\pi/2$,
respectively.  Thus their clockwise lifts have the
relations in \eqref{eq:q4-unit-tangent-central-presentation}, and
\eqref{eq:q4-four-strand-central-lift} respects both relations because
$\Delta_4^{-2}=\delta_4^{-4}=z_4^{-1}$.

The resulting map to the pullback in
\eqref{eq:q4-four-strand-pullback-extension} is an isomorphism on the central
kernel, sending $f^k$ to $z_4^k$.  Given a pair $(\gamma,b)$ in the pullback,
choose a lift $\widetilde\gamma$.  The braids $b$ and the image of
$\widetilde\gamma$ differ by $z_4^k$ for some $k$, so multiplication by $f^k$
gives a preimage.  If an element maps to the identity, its $\Gamma_4$ projection
is trivial and it is therefore $f^k$ for some $k$; the image $z_4^k$ is trivial
only when $k=0$.  This proves injectivity as well as surjectivity, hence the
extension isomorphism.

Finally, the braid exponent of $z_4$ is twelve.  The discriminant map on the
space of square-free monic quartics sends each positive Artin generator to
winding one.  The inverse full twist therefore has discriminant winding
$-12$, proving \eqref{eq:q4-four-strand-normalized-discriminant-period}.
\end{proof}

\begin{corollary}[Central history residue and LHS transgression]
\label{cor:q4-central-history-transgression}
Let $\mathfrak F_4=\langle\mathbf j,\mathbf r\rangle$ be the free reversible
word group, let
\begin{equation}
 \rho:\mathfrak F_4\longrightarrow\Gamma_4,
 \qquad
 \mathbf j\longmapsto J,
 \qquad
 \mathbf r\longmapsto R,
 \qquad
 N=\ker\rho,
 \label{eq:q4-free-word-operator-extension}
\end{equation}
and supply the coefficient-path homomorphism
\begin{equation}
 \widehat\Psi_4:\mathfrak F_4\longrightarrow B_4,
 \qquad
 \mathbf j\longmapsto\Delta_4^{-1},
 \qquad
 \mathbf r\longmapsto\delta_4^{-1}.
 \label{eq:q4-supplied-four-root-path-map}
\end{equation}
There is a unique integral character $\ell:N\to\Z$ such that
\begin{equation}
 \widehat\Psi_4(n)=z_4^{\ell(n)}.
 \label{eq:q4-central-history-character}
\end{equation}
It is $\mathfrak F_4$-conjugation invariant, hence represents a class in
$H^1(N;\Z)^{\Gamma_4}$, and satisfies
\begin{equation}
 \ell(\mathbf j^2)=\ell(\mathbf r^4)=-1,
 \qquad
 \ell(n)=\frac1{12}\operatorname{exp}(\widehat\Psi_4(n))
 =\frac1{12}\frac1{2\pi i}\int_{P_n}\dd\log\operatorname{disc}P_n.
 \label{eq:q4-central-history-period}
\end{equation}
Here $P_n$ is a square-free monic quartic coefficient loop realizing
$\widehat\Psi_4(n)$.  If
$\widetilde\rho:\mathfrak F_4\to\widetilde\Gamma_4$ sends the two generators
to the selected clockwise elliptic lifts and
\begin{equation}
 \widetilde\rho(n)=f^{\operatorname{rot}_{\mathrm{UT}}(n)}
 \qquad(n\in N),
 \label{eq:q4-unit-tangent-relation-rotation}
\end{equation}
then the same integer has the cross-layer readings
\begin{equation}
 \operatorname{rot}_{\mathrm{UT}}(n)
 =\ell(n)
 =\frac1{12}\operatorname{exp}(\widehat\Psi_4(n))
 =\frac1{12}\frac1{2\pi i}\int_{P_n}\dd\log\operatorname{disc}P_n.
 \label{eq:q4-central-residue-cross-layer-identity}
\end{equation}
The pushout of
$1\to N\to\mathfrak F_4\to\Gamma_4\to1$ along $\ell$ is canonically
isomorphic to the unit-tangent extension
$1\to\Z\to\widetilde\Gamma_4\to\Gamma_4\to1$.  Equivalently, with the
Lyndon--Hochschild--Serre convention in which transgression is the class of this
pushout,
\begin{equation}
 d_2[\ell]=e_{\mathrm{UT}}=\psi^*e_{B_4}
 \quad\text{in}\quad H^2(\Gamma_4;\Z).
 \label{eq:q4-central-history-transgression}
\end{equation}
Here $e_{B_4}$ is the class of
$1\to\langle z_4\rangle\to B_4\to\overline B_4\to1$ under the
identification $\langle z_4\rangle\cong\Z$.
With the opposite convention for the LHS differential, the left side is written
$-d_2[\ell]$.  Reversing the chosen generator of all three central copies of
$\Z$ changes the three displayed classes simultaneously.
\end{corollary}

\begin{proof}
The quotient of \eqref{eq:q4-supplied-four-root-path-map} by
$\langle z_4\rangle$ is $\psi\circ\rho$.  Hence $N$ maps into
$\langle z_4\rangle$, which defines the unique exponent $\ell$.  Centrality of
$z_4$ makes $\ell$ invariant under conjugation.  The relator values follow from
$\Delta_4^{-2}=\delta_4^{-4}=z_4^{-1}$, and the general period formula follows
from $\operatorname{exp}(z_4)=12$.

Write the pushout as the quotient of $\Z\times\mathfrak F_4$ by the normal
relations $(-\ell(n),n)=1$ for $n\in N$.  The map
\begin{equation}
 (k,w)\longmapsto\bigl(\rho(w),z_4^k\widehat\Psi_4(w)\bigr)
 \label{eq:q4-central-history-pushout-map}
\end{equation}
is well defined and identifies it with
$\Gamma_4\mathbin{\times}_{\overline B_4}B_4$.  Given $(\gamma,b)$ in the
pullback, choose $w\in\mathfrak F_4$ with $\rho(w)=\gamma$.  Since
$b\widehat\Psi_4(w)^{-1}=z_4^k$ for a unique $k$, the class of $(k,w)$ maps to
$(\gamma,b)$.  The kernel relation is exactly $k=-\ell(w)$ for $w\in N$.
Proposition~\ref{prop:q4-four-strand-central-calibration}
identifies this pullback with $\widetilde\Gamma_4$.  Under the pullback
isomorphism, the two central generators are identified, so
$\operatorname{rot}_{\mathrm{UT}}(n)=\ell(n)$ and
\eqref{eq:q4-central-residue-cross-layer-identity} follows.  The five-term sequence for
the Lyndon--Hochschild--Serre spectral sequence identifies the transgression of
an invariant kernel character with its pushout extension class
\cite[Chapter~VII]{Brown1982CohomologyGroups}, proving
\eqref{eq:q4-central-history-transgression}.
\end{proof}

No injectivity of $\psi$ is needed here.  More importantly, the proposition
and corollary start by supplying a standard root configuration and the paths
$\Delta_4^{-1},\delta_4^{-1}$.  It proves equality of two central extension
classes, not that the q=4 zero dessin or every marked arithmetic history forces
those coefficient paths.  The positive full twist $z_4$ closes to $T(4,4)$,
whereas the supplied inverse relator closes to its mirror $T(4,-4)$.  Both have
four components, and neither is the two-component cusp link $T(2,4)$ of
Theorem~\ref{thm:q4-sign-cover-torus-link}.

The character $\ell$ is a new integral relation residue, not Paper~I's affine
ACS torsion.  It is defined here for a one-object reversible word group.  A typed
history category would require compatible characters on the isotropy kernels of
its objects before an analogous exact sequence or transgression could be stated.

There is nevertheless a direct information-filtration link to Paper~I.  Its
determinant-one calculation, with
$R_{\mathrm{mat}}:=\widetilde J_{\mathrm{mat}}\widetilde T$, gives
$\widetilde J_{\mathrm{mat}}^{\,2}=R_{\mathrm{mat}}^{\,4}=-I$, while the
projective quotient sends both scalars to the identity.  In the covering tower
from the unit-tangent central extension through
$\operatorname{SL}_2(\R)$ to $\operatorname{PSL}_2(\R)$, the positive fiber
generator $f$ maps to $-I$.  Thus, on these central residues, for every
$n\in N$ the three levels retain
\begin{equation}
 z_4^{\ell(n)}
 \longmapsto
 (-I)^{\ell(n)}
 \longmapsto
 1,
 \qquad
 \Z\xrightarrow{\bmod 2}\Z/2\Z\longrightarrow0.
 \label{eq:q4-central-residue-information-filtration}
\end{equation}
The quartic coefficient path and the unit-tangent lift retain the full integer;
the determinant-one matrix retains its parity; projectivization forgets both.
This is an exact process-residue filtration, but it still does not identify
$\ell$ with the continuous affine torsion of Paper~I.

\begin{problem}[Mixed slice--tube naturality]
\label{prob:mixed-slice-tube-naturality}
Construct a typed history family whose slice evaluations produce the coefficient
path of a polynomial $P(z,t)$, and for which history equivalence acts by the
declared toric and divisor gauges.  Determine whether a slice-level arithmetic
torsion class has a well-typed monodromy variation or central-extension lift whose
horizontal period equals the discriminant period of $P$.  The supplied family
$z^2-t^m$, the q=4 toric normal form, and
Proposition~\ref{prop:q4-four-strand-central-calibration} together with
Corollary~\ref{cor:q4-central-history-transgression} prove geometric,
horizontal, central-extension, and relation-residue identities, but not this
unrestricted history-naturality statement.

For an invertible marked-word group $F$ with operator quotient
$1\to N\to F\xrightarrow{\rho}\Gamma_4\to1$, a precise two-sheet target is a
commutative diagram of exact rows
\begin{equation}
\begin{array}{ccccccccc}
1&\longrightarrow&N&\longrightarrow&F&\xrightarrow{\rho}&\Gamma_4&\longrightarrow&1\\
 &&\downarrow&&\downarrow\widehat\Psi&&\downarrow\chi&&\\
1&\longrightarrow&P_2&\longrightarrow&B_2&\longrightarrow&S_2&\longrightarrow&1.
\end{array}
\label{eq:history-to-pure-braid-target}
\end{equation}
The image in $P_2$ of two histories with the same operator would then be a
well-typed horizontal kernel residue, not a Serre transgression.  This target
tests the sign-cover cusp thread.

The four-sheet unit-tangent target is different.  It asks instead for
\begin{equation}
\begin{array}{ccccccccc}
1&\longrightarrow&N&\longrightarrow&F&\xrightarrow{\rho}&\Gamma_4&\longrightarrow&1\\
 &&\downarrow&&\downarrow\widehat\Psi_4&&\downarrow\psi&&\\
1&\longrightarrow&\langle z_4\rangle&\longrightarrow&B_4&\longrightarrow&
\overline B_4&\longrightarrow&1.
\end{array}
\label{eq:history-to-four-braid-target}
\end{equation}
Proposition~\ref{prop:q4-four-strand-central-calibration} constructs this diagram
for the free operator words after the two generator paths are supplied: the
relations $J^2$ and $R^4$ both map to $z_4^{-1}$.  Extending it naturally to the
typed arithmetic histories, their domains, and their declared equivalences is
the open part.  The two-strand target tests the sign-cover cusp thread; the
four-strand target tests the unit-tangent Euler class.  They are different
functors and need not agree.
\end{problem}

The neutral history $\omega_4=(JT)^4$ is a decisive test.  The terminal divisor
of Section~\ref{sec:history-divisor-naturality} sends it to the same endpoint
object as the empty history.  A successful horizontal refinement would have to
produce a coefficient or framed-configuration loop from the full marked
history, prove its gauge independence, and then determine whether that loop is
$\sigma_1^4$ in the two-sheet normalization and $z_4^{-1}$ in the four-sheet
normalization.  Merely attaching either already known torus-link model after
projective evaluation would not pass this test.
